\documentclass[12pt,a4paper]{article}

\usepackage[margin=1in]{geometry}

\title{CPP\_STL Project Note}
\author{C++ STL and Algorithm Practice Workspace}
\date{\today}

\begin{document}

\maketitle
\tableofcontents
\newpage

\section{Project Overview}

The \texttt{CPP\_STL} project is a hands-on C++ learning workspace. Its goal is
to practise the Standard Template Library (STL), fundamental algorithms,
recursion, sorting, and competitive-programming style problem solving. The
repository currently contains 40 C++ source files. These files are mostly small,
independent exercises rather than one combined application.

The project is useful for learning how to select a data structure, translate an
algorithm into C++, read input efficiently, and reason about time complexity.

\section{Tools and Environment}

\begin{itemize}
    \item Language: C++
    \item Typical compiler: \texttt{g++} with the C++17 standard
    \item Editor: VS Code or VSCodium
    \item Documentation build: LaTeX Workshop with \texttt{latexmk}
    \item Main library: the C++ Standard Template Library (STL)
\end{itemize}

For a normal C++ source file, a useful compile command is:

\begin{verbatim}
g++ -std=c++17 -Wall -Wextra -Wpedantic filename.cpp -o program
./program
\end{verbatim}

\section{Topics Covered}

\subsection{Input, Output, and Arrays}

Several files demonstrate basic input/output, arrays, and vectors. Examples
include \texttt{first.cpp}, \texttt{djsn.cpp}, \texttt{dguvhj.cpp},
\texttt{jcvnbk.cpp}, and \texttt{stl.cpp}. The key ideas are reading values
with \texttt{cin}, printing with \texttt{cout}, traversing elements with loops,
and using \texttt{vector} when the size should be managed safely at runtime.

\subsection{STL Containers}

The project explores common STL containers and their intended use:

\begin{itemize}
    \item \texttt{vector}: dynamic contiguous storage and indexed access.
    \item \texttt{queue}: first-in, first-out processing.
    \item \texttt{stack}: last-in, first-out processing.
    \item \texttt{set} and \texttt{multiset}: ordered unique or repeated values.
    \item \texttt{map}: ordered key--value storage.
    \item \texttt{priority\_queue}: repeated access to the highest-priority item.
\end{itemize}

Files such as \texttt{sort\_a2z.cpp}, \texttt{pblm\_queue.cpp},
\texttt{priotrity\_quque.cpp}, \texttt{custom\_PQ.cpp}, and
\texttt{custom\_PQ\_1.cpp} focus on container behaviour and priority queues.

\subsection{Sorting and Custom Comparison}

Sorting is a recurring theme. The standard form is:

\begin{verbatim}
sort(values.begin(), values.end());
\end{verbatim}

For a special order, C++ accepts a comparator. A comparator determines which
element has higher priority or should appear first. This is especially important
for \texttt{priority\_queue}, where the ordering rules can be different from a
normal sorted array.

\subsection{Recursion}

The files \texttt{recursion.cpp}, \texttt{power.cpp}, \texttt{power\_extenc.cpp},
and \texttt{hjghx.cpp} practise recursive thinking. Typical examples include
factorial, exponentiation, and checking whether a target value can be formed
from a collection of coin values. Every recursive algorithm needs a correct
base case and a step that reduces the problem size.

\section{Core Algorithm: Custom Interval Array Construction}

The most complete algorithmic implementation is \texttt{d\_custom\_array.cpp}.
It constructs an array by repeatedly selecting the longest still-empty interval,
placing the next value at its middle, and splitting that interval into two
smaller intervals.

\subsection{Algorithm Steps}

\begin{enumerate}
    \item Start with the interval \texttt{[1, n]}.
    \item Store active intervals in a \texttt{priority\_queue}.
    \item Give higher priority to a longer interval.
    \item If two intervals have equal length, choose the leftmost one.
    \item Place the next number at the middle position. For an even-length
          interval, the implementation uses the left of the two middle cells.
    \item Push the non-empty left and right subintervals back into the queue.
    \item Continue until every array position has received a value.
\end{enumerate}

The comparator captures the selection rule:

\begin{verbatim}
if (len1 != len2) return len1 < len2;
return x.first > y.first;
\end{verbatim}

This makes a larger interval appear at the top of the priority queue. If the
lengths match, the interval with the smaller left endpoint is selected first.

\subsection{Complexity}

There are \(n\) placements. Each interval insertion or removal from the
priority queue costs \(O(\log n)\). Therefore, the overall complexity is
\(O(n \log n)\), while the auxiliary storage is \(O(n)\).

\section{Important Source Files}

\begin{center}
\begin{tabular}{|l|l|}
\hline
\textbf{File} & \textbf{Main learning objective} \\
\hline
\texttt{d\_custom\_array.cpp} & Priority queue, comparator, interval splitting \\
\hline
\texttt{stl.cpp} & Basic string input and string size \\
\hline
\texttt{sort\_a2z.cpp} & Sorting, reversing, multiset, priority queue \\
\hline
\texttt{custom\_PQ.cpp} & Priority queue experiments \\
\hline
\texttt{custom\_PQ\_1.cpp} & Custom ordering with comparators \\
\hline
\texttt{pblm\_queue.cpp} & Stream processing with a priority queue \\
\hline
\texttt{recursion.cpp} & Recursive factorial \\
\hline
\texttt{power.cpp} & Iterative exponentiation \\
\hline
\texttt{power\_extenc.cpp} & Recursive exponentiation \\
\hline
\texttt{hjghx.cpp} & Recursive target/coin search \\
\hline
\end{tabular}
\end{center}

Other practice files include \texttt{fibonbaci.cpp}, \texttt{follw.cpp},
\texttt{probem1.cpp}, \texttt{problem2.cpp}, \texttt{problem3.cpp},
\texttt{problem4.cpp}, \texttt{problem5.cpp}, \texttt{problem\_3.cpp},
\texttt{problen4.cpp}, \texttt{tgt.cpp}, and several short experimental files.

\section{Current Project Structure}

\begin{verbatim}
CPP_STL/
|-- d_custom_array.cpp       # interval-priority algorithm
|-- custom_PQ.cpp            # priority queue practice
|-- custom_PQ_1.cpp          # custom comparator practice
|-- sort_a2z.cpp             # sorting and STL experiments
|-- recursion.cpp            # recursion example
|-- power.cpp                # iterative power
|-- power_extenc.cpp         # recursive power
|-- pblm_queue.cpp           # priority queue problem practice
|-- latex_sample.tex         # minimal LaTeX example
`-- CPP_STL_Project_Note.tex # this report
\end{verbatim}

\section{Code Quality Notes}

This is an educational workspace, so some files are incomplete drafts or
experiments. Before submitting a solution or using a file as a final program,
the following checks are recommended:

\begin{itemize}
    \item Keep one \texttt{main()} function in each compilable source file.
    \item Prefer \texttt{vector<int>} over variable-length arrays in standard C++.
    \item Compile with warnings enabled and fix every warning.
    \item Verify the exact input format, especially whether the first number is
          a test-case count.
    \item Test boundary cases such as \(n=1\), empty input where applicable,
          repeated values, and even-length intervals.
\end{itemize}

\section{Conclusion}

The CPP\_STL project provides a practical foundation for C++ programming and
competitive programming. Its strongest learning themes are STL containers,
priority queues, comparators, sorting, recursion, and interval-based problem
solving. The next natural step is to organise finished solutions into named
folders, keep one problem statement with each solution, and add test cases for
each completed program.

\end{document}
